\section{The full support}
\label{sec:fullsupport}

The full support $S=\{1,\dots,k-1\}$ never carries a $\sigma$-bump. Indeed,
(B1) forces $Q=\{\sigma\}$, and then at $i=k-1$ every position $i+j$ with
$j\in S$ lies in $\{k,\dots,2k-2\}$, where the only position that can lie outside $P$ is
$k+\sigma$, which (B2) excludes. So the full support has to be treated on its
own. It is also the classical case: the full-support recursion, in which every
term of the window appears in the numerator, was settled, even for complex
coefficients, by \citet{CsornyeiLaczkovich2001}. They also passed to its
piecewise-linear shadow and showed, by exhibiting a non-periodic real solution
at every order other than these, that it is periodic only at order $2$ and,
with $c_0>0$, at order $3$ (their Theorems~9--11). The periods of the shadow at order $4$ with $c_0>0$ are determined by
\citet{LineroBasNievesRoldan2021}, and, with $c_0>0$, every solution with monotone initial
data has period $3k-1$ \citep{GolombCallan1992}. The carrier $Q=\{1\}$ satisfies (B1) with $\sigma=1$ but fails (B2), so the bump
is not a tangent solution along it; but every window of the full support
meets this carrier, so \autoref{lem:trop-carrier} applies, and here I give a
different unbounded tangent solution along its orbit, written down in closed
form and checked by hand.

\begin{theorem}
\label{thm:fullsupport}
Let $k\ge4$. No equation \eqref{eq:familya} with full support
$S=\{1,\dots,k-1\}$ is globally periodic, for either value of $\theta$ and for any
choice of positive coefficients.
\end{theorem}

The bound $k\ge4$ is sharp: at $k=2$ and $k=3$ the full supports $\{1\}$ and
$\{1,2\}$ carry the period-$6$, Lyness and period-$8$ equations of
\eqref{eq:fivefamilies}.

\subsection{The carrier}
\label{sec:fs-carrier}

Throughout this section $k\ge4$ and $S=\{1,\dots,k-1\}$. Let $P\subseteq\Z/2k$ be
the anti-periodic set coded by $Q=\{1\}$ (\autoref{def:bump}), that is
\[
P=\{1,\,k\}\cup\{k+2,k+3,\dots,2k-1\},
\]
and let $E$ be its orbit. Write every $n\ge0$ as $n=\ell k+i$ with $\ell\ge0$ and
$0\le i<k$, and call $\ell$ the \textit{block} and $i$ the \textit{column} of $n$.

\begin{lemma}
\label{lem:fs-carrier}
The orbit $E$ is a solution of the min-plus recursion for both values of $\theta$,
with $m_n(E)=-1$, and its attaining sets are
\[
A_{\ell k+i}(E)=\begin{cases}
\{1\}, & \ell\ \text{even},\ i=0,\\
\{k-i\}\cup\{k+2-i,\dots,k-1\}, & \ell\ \text{even},\ i\ge1,\\
\{2,\dots,k-1\}, & \ell\ \text{odd},\ i=0,\\
\{1,\dots,k-1-i\}\cup\bigl(\{k+1-i\}\cap S\bigr), & \ell\ \text{odd},\ i\ge1.
\end{cases}
\]
\end{lemma}

\begin{proof}
Let $1\le j\le k-1$. If $\ell$ is even, $n+j\equiv i+j\pmod{2k}$, with
$1\le i+j\le2k-2$. When $i+j<k$ this lies in $P$ only if $i+j=1$, that is $i=0$
and $j=1$; when $i+j\ge k$ it lies in $P$ exactly when $i+j=k$ or $i+j\ge k+2$.
If $\ell$ is odd, $n+j\equiv k+i+j\pmod{2k}$. When $i+j<k$ this lies in $P$
exactly when $i+j\ne1$, which for $j\ge1$ means $i\ge1$ or $j\ge2$; when $i+j\ge k$
the residue is $i+j-k\in\{0,\dots,k-2\}$, which lies in $P$ only if $i+j=k+1$.
This gives the four sets $\{j\in S:n+j\in P\}$ listed, and each of them is nonempty
(it contains $1$, $k-i$, $2$, and $1$ or $2$ respectively). So
\autoref{lem:trop-carrier} applies.
\end{proof}

In words, a tap $j$ at $n=\ell k+i$ reads column $i+j$ of block $\ell$ if
$i+j<k$, and column $i+j-k$ of block $\ell+1$ otherwise. So in an even block the
attaining taps read column $1$ of the same block (if $i=0$), or column $0$ and
columns $2,\dots,i-1$ of the next block (if $i\ge1$). In an odd block they read
columns $i+1,\dots,k-1$ of the same block (columns $2,\dots,k-1$ if $i=0$), and
also column $1$ of the next block when $i\ge2$. For example, at $k=5$ the
columns read by the attaining taps are as follows, a prime marking a column of
the next block:
\begin{center}
\small
\begin{tabular}{l|ccccc}
column $i$ of $n$ & $0$ & $1$ & $2$ & $3$ & $4$\\\hline
$\ell$ even & $1$ & $0'$ & $0'$ & $0',2'$ & $0',2',3'$\\
$\ell$ odd & $2,3,4$ & $2,3,4$ & $3,4,1'$ & $4,1'$ & $1'$
\end{tabular}
\end{center}
Here the middle columns are $2$ and $3$.

\subsection{The ray}

Write the block as $\ell=8m+r$ with $m\ge0$ and $0\le r<8$, so that
$n=8km+rk+i$. Define $d_n$ by \autoref{tab:fs-ray}: its value depends on $r$, on
$m$, and on the column only through the four classes $i=0$, $i=1$,
$2\le i\le k-2$ (the \textit{middle} columns) and $i=k-1$, which are distinct and
nonempty because $k\ge4$. Write $a_r$, $b_r$, $c_r$ and $z_r$ for the entries of
row $r$ in these four classes. Only the middle entries depend on $m$: they are
$f-m$ in even blocks and $f+m$ in odd blocks, for a constant $f$ given by the row.
Block $8$, which is block $0$ of the next group, has $a_8=a_0$, $b_8=b_0$,
$z_8=z_0$ and $c_8=-(m+1)$. So $d$ is a ray (\autoref{def:trop-ray}) with $q=8k$: passing from $m$ to $m+1$ changes only the
middle entries, by $-1$ in even blocks and by $+1$ in odd ones.

\begin{table}
\centering
\small
\setlength{\tabcolsep}{4pt}
\begin{tabular}{c|cccc|cccc}
$r$ & $a_r$ & $b_r$ & $c_r$ & $z_r$ &
$a_r+a_{r+1}$ & $b_r+b_{r+1}$ & $c_r+c_{r+1}$ & $z_r+z_{r+1}$\\
& $(i=0)$ & $(i=1)$ & $(2\le i\le k-2)$ & $(i=k-1)$ &&&&\\
\hline
$0$ & $-1$ & $0$ & $-m$ & $0$ & $0$ & $1$ & $1$ & $1$\\
$1$ & $1$ & $1$ & $m+1$ & $1$ & $1$ & $1$ & $0$ & $0$\\
$2$ & $0$ & $0$ & $-m-1$ & $-1$ & $0$ & $0$ & $0$ & $0$\\
$3$ & $0$ & $0$ & $m+1$ & $1$ & $1$ & $1$ & $1$ & $1$\\
$4$ & $1$ & $1$ & $-m$ & $0$ & $1$ & $0$ & $0$ & $0$\\
$5$ & $0$ & $-1$ & $m$ & $0$ & $0$ & $0$ & $0$ & $1$\\
$6$ & $0$ & $1$ & $-m$ & $1$ & $1$ & $1$ & $1$ & $1$\\
$7$ & $1$ & $0$ & $m+1$ & $0$ & $0$ & $0$ & $0$ & $0$\\
\end{tabular}
\caption{The ray of \autoref{lem:fs-ray}. Left: the value of $d_n$ at
$n=8km+rk+i$, by column class. Right: the left-hand side $d_{n+k}+d_n$ of the
tangent recursion, by column class; for $r=7$ the next block is block $0$ with $m$
replaced by $m+1$.}
\label{tab:fs-ray}
\end{table}

\begin{lemma}
\label{lem:fs-ray}
The sequence $d$ of \autoref{tab:fs-ray} is a tangent solution along $E$, for both
values of $\theta$, and $d_{8km+2k+2}=-1-m$ for every $m\ge0$.
\end{lemma}

\begin{proof}
Let $n=\ell k+i$ with $\ell=8m+r$. The left-hand side $d_{n+k}+d_n$ is the sum of
the entries of rows $r$ and $r+1$ in the class of column $i$. In the middle class
the terms in $m$ cancel, and the four sums are the right-hand half of
\autoref{tab:fs-ray}.

Now the right-hand side, the minimum of $d_{n+j}$ over $j\in A_n(E)$. Note first
that in every odd row $c_r\ge a_r$ and $c_r\ge z_r$, since $m\ge0$. By the
description after \autoref{lem:fs-carrier}, the middle columns are read only
from odd blocks, and always alongside a column whose entry is no larger: in an
even block $\ell$ at $i\ge3$ they are read from block $\ell+1$ together with its
column $0$, and in an odd block at $i\le k-3$ they are read together with column
$k-1$ of the same block. So the middle reads never lower the minimum and can be ignored,
and the right-hand side is:
\begin{itemize}
\item for $r$ even: $b_r$ if $i=0$, and $a_{r+1}$ if $i\ge1$;
\item for $r$ odd: $z_r$ if $i\in\{0,1\}$, $\min(z_r,b_{r+1})$ if
$2\le i\le k-2$, and $b_{r+1}$ if $i=k-1$.
\end{itemize}
The tangent recursion \eqref{eq:trop-tangent} therefore reduces to the identities
\[
\begin{aligned}
&r\ \text{even:}&& a_r+a_{r+1}=b_r,\quad
b_r+b_{r+1}=c_r+c_{r+1}=z_r+z_{r+1}=a_{r+1};\\
&r\ \text{odd:}&& a_r+a_{r+1}=b_r+b_{r+1}=z_r,\quad
c_r+c_{r+1}=\min(z_r,b_{r+1}),\quad z_r+z_{r+1}=b_{r+1}.
\end{aligned}
\]
These are read off \autoref{tab:fs-ray}. For the even rows $r=0,2,4,6$ the values
$(b_r;a_{r+1})$ are $(0;1)$, $(0;0)$, $(1;0)$ and $(1;1)$, and for the odd rows
$r=1,3,5,7$ the values $(z_r;\min(z_r,b_{r+1});b_{r+1})$ are $(1;0;0)$,
$(1;1;1)$, $(0;0;1)$ and $(0;0;0)$; each matches the corresponding sums. By
\autoref{lem:fs-carrier} the attaining sets, and hence this computation, are the
same for both $\theta$.

Finally, column $2$ is a middle column because $k\ge4$, so
$d_{8km+2k+2}=c_2=-1-m$.
\end{proof}

\begin{proof}[Proof of \autoref{thm:fullsupport}]
By \autoref{lem:fs-carrier} the orbit $E$ is an integer solution of the min-plus
recursion, and by \autoref{lem:fs-ray} it carries an unbounded integer tangent
solution, for both values of $\theta$. \autoref{thm:tangent} applies.
\end{proof}
